% ****** Start of file apssamp.tex ******
%
%   This file is part of the APS files in the REVTeX 4.1 distribution.
%   Version 4.1r of REVTeX, August 2010
%
%   Copyright (c) 2009, 2010 The American Physical Society.
%
%   See the REVTeX 4 README file for restrictions and more information.
%
% TeX'ing this file requires that you have AMS-LaTeX 2.0 installed
% as well as the rest of the prerequisites for REVTeX 4.1
%
% See the REVTeX 4 README file
% It also requires running BibTeX. The commands are as follows:
%
%  1)  latex apssamp.tex
%  2)  bibtex apssamp
%  3)  latex apssamp.tex
%  4)  latex apssamp.tex
%


\documentclass[%
 reprint,
superscriptaddress,
longbibliography,
%groupedaddress,
%unsortedaddress,
%runinaddress,
%frontmatterverbose, 
%preprint,
%showpacs,preprintnumbers,
%nofootinbib,
%nobibnotes,
%bibnotes,
 amsmath,amssymb,
 prl,
 citeautoscript,
%pra,
%prb,
%rmp,
%prstab,
%prstper,
%floatfix,
]{revtex4-1}

\usepackage{graphicx}% Include figure files
\usepackage{dcolumn}% Align table columns on decimal point
\usepackage{bm}% bold math
\usepackage{amsmath}
\usepackage{amssymb}
\usepackage{amsthm}
\usepackage{units}
\usepackage{gensymb}
\usepackage{xcolor}
\newcommand{\tN}{t_\mathrm{N}}
\newcommand{\kB}{k_\mathrm{B}}


\DeclareMathOperator{\sgn}{sgn}
\DeclareMathOperator{\eff}{eff}

\begin{document}


\title{Emergent Kelvin waves in chiral active matter}

\author{Anthony R. Poggioli}
\email{arpoggioli@berkeley.edu}
\affiliation{Department of Chemistry, University of California, Berkeley}
\affiliation{Kavli Energy NanoScience Institute, Berkeley, California}

\author{David T. Limmer}
\email{dlimmer@berkeley.edu}
\affiliation{Department of Chemistry, University of California, Berkeley}
\affiliation{Kavli Energy NanoScience Institute, Berkeley, California}
\affiliation{Materials Science Division, Lawrence Berkeley National Laboratory}
\affiliation{Chemical Science Division, Lawrence Berkeley National Laboratory}



\date{\today}

\maketitle

%\noindent {\bf The phenomenological equations of hydrodynamics describe emergent phenomena in many body systems. Accordingly, their derivation from first principles is of primary importance in understanding collective phenomena in complex materials. However, traditional frameworks for the derivation of hydrodynamic equations rely on the assumption that the quiescent state of the system is one of thermodynamic equilibrium. We present a general framework for deriving hydrodynamic equations for systems far from equilibrium and apply our method to a chiral active fluid, in which both parity and time-reversal symmetries are broken. Our equations successfully rationalize the emergence of a spontaneous boundary current in the confined fluid, allowing us to extract estimates of transport coefficients that we favorably compare to forced runs. Our hydrodynamic solution reveals that the boundary current is exactly analogous to a quasigeostrophic coastal current, a well known phenomenon in coastal oceanography. Such currents are conjugate to a class of chiral waves called coastal Kelvin waves. Motivated by this analogy, we demonstrate that an acoustic chiral Kelvin wave mode also exists in confined chiral active matter in the absence of an imposed rotation, originating from the spontaneous emergence of a Coriolis-like parameter in the bound modes of a chiral fluid.}

\noindent {\bf The phenomenological equations of hydrodynamics describe emergent behavior in many body systems. %Accordingly, their derivation from first principles is of primary importance in understanding collective phenomena in complex materials. 
Their forms and the associated phenomena are well established when the quiescent state of the system is one of thermodynamic equilibrium, yet away from equilibrium relatively little is firmly established. Here, we deduce directly from first principles the hydrodynamic equations for a system far from equilibrium, a chiral active fluid in which both parity and time-reversal symmetries are broken. With our theory, we rationalize the emergence of a spontaneous boundary current in the confined fluid, a feature forbidden at equilibrium, which allows us to extract estimates of transport coefficients that we favorably compare to forced flows. The hydrodynamic solution reveals that the boundary current is analogous to a quasigeostrophic coastal current, a well known phenomenon in oceanography. Such currents are conjugate to a class of chiral waves called Kelvin waves. Motivated by this analogy, we demonstrate that an acoustic chiral Kelvin wave mode also exists in confined chiral active matter in the absence of an imposed rotation, originating from the spontaneous emergence of a Coriolis-like parameter in the bound modes of a chiral fluid.}


Chiral active fluids are composed of actively rotating constituents, and this activity breaks both time-reversal and parity symmetries. This leads to the emergence of an exotic transport coefficient called the Hall or odd viscosity, $\eta_O$ \cite{Avron}. Qualitatively, $\eta_O$ has the effect of linking forcing to orthogonal fluid motion \cite{cargo, Hall-like}, and is therefore phenomenologically similar to the Coriolis parameter that appears in oceanic and atmospheric flows \cite{Gill, science_topology, gauge_theory}. We show that this analogy becomes exact when considering the boundary-trapped collective motions in a chiral fluid, and we exploit this analogy to demonstrate the existence of a chiral, boundary-trapped acoustic wave mode, analogous to oceanographic coastally bound Kelvin waves.

\begin{figure*}[t]
	\centerline{\includegraphics[width=17cm]{FIg1_vDL.pdf}}
    \caption{Structure of spontaneous boundary current in a chiral active fluid. a), Velocity profile obtained for ${\rm Pe} = 5$, $\gamma = 0.1 \, \sqrt{m \varepsilon \sigma^{-2}}$. b), Spin anomaly profile for same parameters. c) Velocity and d), spin anomaly profiles collapsed according to theoretical prediction for profile structure. Inset in a is a system snapshot. Gray shaded regions in a and b indicate region beyond hydrodynamic radius. Solid lines are theoretical predictions from Eq. \ref{eqn:v_BC} and \ref{eqn:DL_BC}}
\label{fig:Fig1}
\end{figure*}

In addition to the odd viscosity, broken parity and time-reversal symmetries also lead to the spontaneous formation of a boundary current in confined chiral fluids, with its propagation direction determined by the chirality of the system-- or equivalently, by the sign of the odd viscosity \cite{waves_Kranthi, free_surface, lubensky_etaR6, crystal}. This spontaneous collective current is a consequence of broken detailed balance, characterized by the emergence of microscopic currents, along with the global rotational symmetry of the system, which requires the organization of microscopic currents into a collective, angular momentum carrying, boundary current. Previous studies have treated the boundary current phenomenologically, suggesting that its formation is due to a hypothetical dissipative coupling between the fluid vorticity and rotor frequency, along with the frustration of free spin at the particle boundary \cite{lubensky_etaR6, waves_Kranthi,free_surface,crystal}. This frustration leads to a local suppression of the rotor frequency, and, nominally, to a compensating increase in the fluid vorticity. However, inhomogeneities in hydrodynamic densities within a few particle diameters of an interface are expected in all hydrodynamic systems, and the frustration of spin in this microscopic contact layer need not have any influence on the hydrodynamic structure of the boundary current beyond this layer. We show that the boundary current forms in the absence of spin-vorticity coupling, deduced directly from a first principles evaluation of the emergent hydrodynamics of the system. Its formation is most clearly understood as a consequence of angular momentum conservation (SI). Injection of angular momentum through particle activity is balanced globally by substrate friction. However, frustration of free spin in the system at finite density due to inter-particle interactions leads to a reduction of the total spin and therefore to a compensating increase in orbital angular momentum. This finite orbital angular momentum is transported by the boundary current.

We consider a minimal model of a two-dimensional chiral active fluid of dimers \cite{Kranthi_hydro_eqns, Cory_MD, my_PRL} composed of harmonically bonded monomers of diameter $\sigma$, in which all non-bonded monomers interact via the short-range, repulsive Weeks-Chandler-Andersen potential with characteristic energy $\epsilon$ \cite{WCA}. In our initial study, the dimers are confined by a frictionless circular wall of radius $R = 75 \, \sigma$ (Fig. \ref{fig:Fig1}). Each monomer is subjected to an active force of magnitude $F_a$ perpendicular to its harmonic bond, such that the dimers experience an active torque of average magnitude $F_a d_0$, where $d_0$ is the equilibrium bond length. The activity is quantified by a P\'{e}clet number ${\rm Pe} \equiv 2 \beta F_a d_0$, where $\beta$ is the inverse temperature times Boltzmann's constant (SI). %Complete simulation details are given in the Supplementary Information (SI) (SI CITATION).

We introduce in the SI a framework for deriving hydrodynamic balance laws in systems arbitrarily far from equilibrium inspired by Ref. \cite{Irving+Kirkwood} and modified for stochastic, rather than Hamiltonian, dynamics \cite{dean, underdamped_dean}. Using this framework, we obtain the hydrodynamic equations directly from our microscopic equations of motion. We find for an incompressible fluid of density $\rho_0$
\begin{equation}
\nabla \cdot {\boldsymbol u} = 0
\label{eqn:incompressible_continuity}
\end{equation}
\begin{equation}
\rho_0 D_t = - \frac{\gamma}{m} \rho_0 {\boldsymbol u} - \nabla p + \left( \eta_S + \eta_R \right) \nabla^2 {\boldsymbol u} + \eta_O \nabla^2 {\boldsymbol u}^*
\label{eqn:incompressible_linmom}
\end{equation}
and
\begin{equation}
\rho_0 D_t \Delta \mathcal{L} = - \frac{\gamma}{m} \rho_0 \Delta \mathcal{L} + 2 \eta_R \omega
\label{eqn:incompressible_spin}
\end{equation}
with ${\boldsymbol u}$ the velocity field, $\Delta \mathcal{L}$ the spin density relative to the constant bulk value in the unforced steady state, $D_t \equiv \partial_t + {\boldsymbol u} \cdot \nabla$ the material derivative, $\omega \equiv \left( \nabla^{3d} \times {\boldsymbol u}  \right) \cdot \hat{\boldsymbol z}$ the vorticity orthogonal to the plane of motion, $\gamma$ a linear substrate friction coefficient, and $m$ the monomer mass. These equations are statements of conservation of mass, linear momentum, and spin angular momentum, respectively. In addition to the ordinary shear viscosity $\eta_S$, two exotic transport coefficients appear in Eqs. \ref{eqn:incompressible_linmom} and \ref{eqn:incompressible_spin}: the rotational viscosity $\eta_R$, which couples the antisymmetric part of the stress tensor to the vorticity and has the effect of enhancing the diffusion of linear momentum, and the odd viscosity $\eta_O$ \cite{Avron, Kranthi_hydro_eqns, Banerjee, Kranthi1, bubble_dynamics_odd_viscosity, Vitellis_pretty_pictures, Liao_PDFmechanics, 3d_odd_viscosity}. The odd viscosity is linked to orthogonal fluid motion in Eq. \ref{eqn:incompressible_linmom} by the dual velocity ${\boldsymbol u}^* \equiv \varepsilon \cdot {\boldsymbol u}$, where $\varepsilon$ is the Levi-Civita tensor. Geometrically, a dual vector ${\boldsymbol a}^*$ is the vector ${\boldsymbol a}$ rotated clockwise by $90^\circ$. The microscopic model underpinning these hydrodynamic equations has been demonstrated to exhibit odd viscosity \cite{Cory_MD, my_PRL}. 

In confinement, the fluid spontaneously generates a boundary current (Fig. \ref{fig:Fig1}). The tangential velocity $v$ and spin anomaly profiles are derived by solving Eqs. \ref{eqn:incompressible_continuity}- \ref{eqn:incompressible_spin} subject to the integral statement of angular momentum conservation (SI)
\begin{equation}
\left\langle Q_{\rm hyd} \right\rangle_0 = \int_0^{R_{\rm hyd}} 2 \pi r {\rm d} r \, \rho_0 r v(r)
\label{eqn:integral_angmom}
\end{equation}
where $r$ is measured from the center of the confinement, $\langle \rangle_0$ indicates a nonequilibrium average in the unforced steady state, and $Q_{\rm hyd}$ is the total orbital angular momentum in the hydrodynamic region of the boundary current, which must be measured in simulation. The hydrodynamic description fails close to the interface where microscopic correlations persist (Figs. \ref{fig:Fig1}a and \ref{fig:Fig1}b) requiring us to define a hydrodynamic radius $R_{\rm hyd}$, which characterizes the region over which a hydrodynamic description is valid \cite{Bocquet_Barrat1994}. The hydrodynamic radius $R_{\rm hyd}$ is taken to coincide with the second density peak \cite{Chen_interface,my_JPCL}.

The solutions for the velocity and spin anomaly profiles $v(r)$ and $\Delta \mathcal{L}(r)$ subject to Eq. \ref{eqn:integral_angmom} are given by
\begin{equation}
v(r) = \frac{\kappa_{\gamma}}{2 \rho_0} \frac{\left\langle Q_{\rm hyd} \right\rangle_0}{\pi R_{\rm hyd}^2} e^{-\kappa_{\gamma} \left( R_{\rm hyd} - r \right)}
\label{eqn:v_BC}
\end{equation}
and
\begin{equation}
\Delta \mathcal{L} (r) = \frac{2 \eta_R}{\eta_S + \eta_R} \kappa_{\gamma}^{-1} v(r)
\label{eqn:DL_BC}
\end{equation}
with the inverse length scale $\kappa_{\gamma} \equiv \sqrt{(\gamma/m) \rho_0 / \left(\eta_S + \eta_R\right)}$. The theoretical profiles depend on only two transport coefficients, $\eta_S$ and $\eta_R$, allowing us to extract estimates of these parameters from fits of the velocity and spin profiles. We measure $\left\langle Q_{\rm hyd} \right\rangle_0 / \pi R_{\rm hyd}^2$ in each simulation, extract $\kappa_{\gamma}$ as the single free parameter in an exponential fit to the velocity profile, and estimate $\eta_R$ from a fit to $\Delta \mathcal{L}$ with $\kappa_{\gamma}^{-1}$ as the imposed decay length. This procedure collapses our data excellently (Figs. \ref{fig:Fig1}c and \ref{fig:Fig1}d), indicating the validity of our theory. This indicates that the flow is incompressible and well described by our first principles hydrodynamic equations.

We test the values of $\eta_S$ and $\eta_R$ extracted from the boundary current simulations against forced runs conducted in a square periodic domain of side length $L = 100 \, \sigma$ in which the chiral active fluid is subjected to a force density ${\boldsymbol f}(y) = (\rho_0 F_0/m) {\rm sin} \left( k y \right)$, with $k = 10 \pi / L$. The anticipated theoretical profiles in this configuration depend on $\eta_S$ and $\eta_R$ and are derived in the SI. Figs. \ref{fig:Fig2}a and b show comparisons of forced simulation results to the theoretical profiles. The agreement is excellent over a range of microscopic parameters, further validating our hydrodynamic theory. This extends the previously observed consistency between linear response relations and forced runs \cite{Kranthi1, Cory_MD} observed for bulk-like homogeneous fluids in periodic simulations to the boundary value problem posed by the boundary current that forms spontaneously when the reference steady state is confined.

\begin{figure}[t]
	\centerline{\includegraphics[width=8.5cm]{Fig2_vDL.pdf}}
    \caption{Forced runs. a) Results are obtained for rotating dimers in a periodic box subjected to a sinusoidal applied force. b) Velocity and spin anomaly profile for $\gamma = 0.1, 0.3$ and $0.5 \, \sqrt{m \varepsilon \sigma^{-2}}$ and ${\rm Pe = 5}$. Theoretical comparisons are plotted with values of $\eta_S$ and $\eta_R$ inferred from boundary current. c) Estimated values of rotational viscosity from boundary current. d) Origin of rotational viscosity from local solid body rotation of fluid.}
\label{fig:Fig2}
\end{figure}

The rotational viscosity in our theory links the antisymmetric part of the stress tensor $\tau$ to the vorticity: $p^* \varepsilon \equiv \left( \tau - \tau^T \right)/2$, and $p^* = p^*_0 + \eta_R \omega$, with $p^*_0$ the bulk value of the antisymmetric stress in the unforced steady state (SI). This is distinct from how $\eta_R$ is typically defined \cite{etaR1, etaR1, etaR3, etaR4, etaR5, lubensky_etaR6}: $p^* = p^*_0 + \eta_R' \left( 2 \Omega - \omega \right)$ with $\Omega = 2m \mathcal{L} / I_0$ the dimer rotation frequency, and $I_0 \equiv md_0^2/2$ the dimer moment of inertia. This spin-vorticity coupling has previously been cited as the origin of boundary current formation \cite{lubensky_etaR6, waves_Kranthi, free_surface, crystal} and assumes a dissipation through friction induced by inter-rotor collisions that tends to result in a spin density $\mathcal{L} = (I_0/2m) (\omega/2)$. Because the repulsive interactions between monomers in our model are central, there is no such frictional dissipation, and this coupling need not be present \cite{etaR1,Kranthi1}.
%Additionally, the assumption that the spin density reduces to a steady state value determined by a rotational frequency $\omega/2$ conflates the actual spin associated with particle rotation with the kinematic `spin-like' quantity $\omega$. In reality, the two are decoupled, with one being associated with the sum of dimer spins in a fluid element and the other with the net center of mass motion of the fluid element.
We allow for the more general relation $p^* = p^*_0 + \eta_R \omega - \phi \Delta \mathcal{L}$, with $\phi$ a hypothetical transport coefficient linking the spin density to the antisymmetric stress. Our results show that $\phi$ does not play a significant role in the dynamics, and the coupling between spin and linear momentum is therefore one-way, with $\omega$ acting as a source for $\Delta \mathcal{L}$ in Eq. \ref{eqn:incompressible_spin}, but $\Delta \mathcal{L}$ playing no role in the linear momentum dynamics (Eq. \ref{eqn:incompressible_linmom}).

The coupling between stress and vorticity would be prohibited by the principle of objectivity in achiral fluids \cite{Kranthi1} but is permitted here. However, it is the breaking of Galilean, rather than parity, symmetry that results in a nonzero value of $\eta_R$. Galilean invariance is broken by a finite value of the substrate friction $\gamma$ and the corresponding introduction of a privileged substrate velocity. The substrate friction induced by the solid body rotation of a fluid element must be compensated by the net torque $p^*$ acting on the fluid element, leading to a finite $\eta_R$. This is corroborated in Figs. \ref{fig:Fig2}c and d, where we see that $\eta_R$ is only weakly dependent on $\rm Pe$, increases monotonically with $\gamma$, and extrapolates to zero as $\gamma$ vanishes.

Our hydrodynamic theory for the spontaneous boundary current reveals that this phenomenon is analogous to a well known phenomenon in coastal oceanography, the formation of a quasigeostrophic coastal current due to the local injection of freshwater into the coastal ocean (SI). The oceanographic coastal current forms to conserve freshwater \cite{Alex_ARFM, Alex_JFM}, analogous to the conservation of angular momentum that results in the formation of the boundary current (Eq. \ref{eqn:integral_angmom}). This analogy reveals that $\eta_O$ results in the spontaneous emergence of a Coriolis-like parameter $\eta_O \kappa_{\gamma}^2 / \rho_0$ proportional to the odd viscosity and the  square of the inverse confinement length $\kappa_{\gamma}^2$.

The direction of propgation of the coastal current is determined by that of coastally bound Kelvin waves, an oceanographic chiral edge mode \cite{continuum_topology, Kelvin_quantum_Hall} conjugate to the steady coastal curent. The chirality of this boundary mode is in accord with the band topographic notion of bulk-boundary correspondence: the presence of the topological order parameter $f$ breaks parity and time-reversal symmetry in the bulk of the fluid, introducing a band gap between bulk low frequency Rossby and high frequency Poincar\'{e} waves and therefore a nontrivial band topology \cite{science_topology}. This nontrivial topology is associated with an integer-valued topological invariant known as the Chern number, and the difference in these numbers across either the topological interface at the equator where $f$ changes sign or at continental boundaries where $f$  vanishes on the continent corresponds to the net number of permitted chiral modes at the boundary or interface \cite{science_topology, fuchs, boundary_waves, topology_soft_matter, Tauber_JFM}. The resultant edge bands exist in the bulk band gap \cite{fuchs}.

There is no imposed Coriolis force in our system, and the bulk band topology is therefore trivial \cite{boundary_waves} (SI). Bulk-boundary correspondence then predicts the absence of a net number of chiral edge modes. However, we are motivated by the spontaneous emergence of a Coriolis-like parameter in the boundary current solution to search for Kelvin-like chiral wave modes in our system. We search for one-dimensional, bound acoustic waves by relaxing the assumption of incompressibility in our hydrodynamic equations (SI) and allowing for linear perturbations $\Delta \rho$ and $\Delta v$ to the density and tangential velocity. This leads to the following system of equations:
\begin{equation}
\partial_t \Delta \rho = - \rho_0 \partial_y \Delta v
\label{eqn:unsteady_continuity}
\end{equation}
\begin{equation}
0 = - c^2 \partial_r^2 \Delta \rho + \eta_O \left( \partial_r^2 + \partial_y^2 \right) \Delta v
\label{eqn:unsteady_xmom}
\end{equation}
and
\begin{equation}
\rho_0 \partial_t \Delta v = - c^2 \partial_y \Delta \rho
\label{eqn:unsteady_ymom}
\end{equation}
with $c \equiv \sqrt{\partial_{\rho} p|_{\rho_0}}$ the bulk speed of sound when $\eta_O = 0$ and $y \equiv r \theta$ the circumferential coordinate. We have neglected frictional terms and coupling of the waves to the boundary current flow (SI). These equations are valid in the limit that the confinement radius $R$ becomes large compared to the wave thickness $\ell_r \equiv \kappa_r^{-1}$. This is the only regime in which boundary waves may exist in the chiral fluid (SI). As waves penetrate further into the bulk, their frequencies are shifted as a function of $r$ by the changing radius of curvature of the wave characteristic and an along-flow odd viscous pressure gradient that depends on the local radius of curvature. We demonstrate in the SI that no solution can exist to these equations when the frequency is radially dependent. This implies that the one-dimensional boundary waves devolve into two-dimensional bulk waves if their thicknesses become comparable to the confinement radius.

Eqs. \ref{eqn:unsteady_continuity} through \ref{eqn:unsteady_ymom} generically admit two solutions with frequencies $\omega_{\pm}^n = \pm S_{\eta_O} ck_y^n$, with $S_{\eta_O}$ the sign of the odd viscosity, and $k_y^n \equiv 2 \pi / \lambda_y = n/R$ the circumferential wave number. Here $n$ is a natural number, and $k_y^n = n/R$ is the periodicity condition in circular confinement. We plot in Fig. \ref{fig:Fig3}a the $e$-folding thicknesses of these modes as a function of circumferential wavelength $\lambda_y$ and normalized by the Rossby radius of deformation $\ell_R \equiv |\eta_O|/\rho_0 c$. The Rossby radius is defined generically as $\ell_R \equiv c/|f|$ in the oceanographic context and sets the thickness of the coastal Kelvin wave \cite{Kundu}. The confinement scale of the $-S_{\eta_O}$ mode diverges quadratically as $\lambda_y \rightarrow \infty$ and therefore must devolve into two dimensional bulk waves beyond some critical wavelength. This mode can only exist when $n \gg (1/2) \sqrt{1 + 2 \kappa_R R}$ (SI). That is, as the system size becomes large compared to the Rossby radius the negatively propagating mode is entirely eliminated, and we are left only with the Kelvin mode. The dispersion of this mode is compared to the bulk dispersion relations in Fig. \ref{fig:Fig3}b. Interestingly, this band fills the bulk band gap everywhere except at $k_y^n = 0$, where all of the frequencies vanish. 

Fig. \ref{fig:Fig3}a reveals that the thickness of the Kelvin mode saturates at the Rossby radius for large wavelengths, analogous to the coastal Kelvin wave. The Rossby radius $\ell_R \equiv |\eta_O|/\rho_0 c$  is proportional to our rotational parameter $\eta_O$, and thus increases in thickness with increasing rotational activity. The Kelvin acoustic mode is illustrated in Fig. \ref{fig:Fig3}c, which microscopically corresponds to alternating rarefaction or compaction of the active particles. Neither this mode nor the oceanographic coastal Kelvin wave are protected in the topological sense. In both cases, this is because their emergence is dependent on the no-flux boundary condition \cite{continuum_topology}. However, these waves are physically protected. As there is only one boundary wave mode with a sense of propagation dictated by the sign of the odd viscosity for system sizes $\kappa_R R \gg 1$, these waves must propagate without backscattering.


\begin{figure}[t]
	\centerline{\includegraphics[width=8.5cm]{Fig3_vDL.pdf}}
    \caption{Chiral Kelvin wave mode. a) Penetration length of Kelvin (red) and divergent (black) wave mode. b) Bulk dispersion relations (blue, gray) and Kelvin edge mode dispersion relation (red). 
    %c) Thickness of Kelvin wave as odd viscosity is varied for fixed transverse wavenumber.
     c) Theoretical structure of boundary-trapped acoustic Kelvin wave. Black dot in b indicate parameter values for wave plotted in c.}
\label{fig:Fig3}
\end{figure}

%We have defined a Rossby radius $\ell_R \equiv |\eta_O|/\rho_0 c$ but have not justified this definition. The oceanographic definition $\ell_R \equiv c/|f|$, leads to the identification of a spontaneously emerging Coriolis parameter $f = \rho_0 c^2 / \eta_O$ in our wave solution. Likewise, inserting our bound wave solution into Eq. \ref{eqn:unsteady_xmom} gives $0 = -c^2 \kappa_R^2 \hat{\Delta \rho} + \eta_O \left( \kappa_R^2 - k_y^2 \right) \hat{\Delta v}$ (SI). Up to the additional dependence on $k_y^n$, this is identical to the corresponding equation for a coastal Kelvin wave with $f = \eta_O \kappa_R^2 / \rho_0 = \rho_0 c^2 / \eta_O$. Notably, the Rossby radius we have identified is proportional to our rotational parameter $\eta_O$, unlike the oceanographic counterpart $\ell_R \propto |f|^{-1}$. We plot $\ell_R$ as a function of $\eta_O$ for fixed $\lambda_y$ in Fig. \ref{fig:Fig3}c, revealing that the boundary mode increases in thickness with increasing rotational activity.


The Kelvin mode discovered here is qualitatively analogous to the chiral free surface mode observed in Ref. \cite{free_surface} for a system of rotating colloids, whose explanation was phenomenological. 
This suggests that the oceanographic analogy here is indeed useful in rationalizing the behavior of confined chiral active fluids. 
The mechanism that underpins this analogy is the spontaneous emergence of a Coriolis-like parameter in boundary trapped motions, and pushing this analogy further will reveal new physics analogous to, for example, two-dimensional instabilities and turbulence in the oceanographic context. This mode was revealed by a first principles derivation of the hydrodynamic balance equations, and this framework allows for the unambiguous derivation of hydrodynamic equations in systems far from equilibrium. Undoubtably, other exotic phenomena are waiting to be uncovered through analogous systematic evaluation of hydrodynamic equations of other classes of matter driven far from equilibrium.

\vspace{0.5cm}
\noindent{\bf Acknowledgments}
This work has been supported by NSF Grant No. CHE-1954580. A. R. P was also supported by the Heising-Simons Fellowship from the Kavli Energy Nanoscience Institute at UC Berkeley and D. T. L acknowledges support from the Alfred P. Sloan Foundation.
\vspace{0.5cm}

%\bibliography{current_ref}

%\section{Methods}
%
%All simulations were conducted in LAMMPS (CITATION). The monomer positions and momenta were evolved according to the underdamped Langevin equation of motion (KRANTHI, TRACER)
%\begin{equation}
%\hat{\boldsymbol p}_i^{\alpha} = - \frac{\gamma}{m} {\boldsymbol p}_i^{\alpha} + {\boldsymbol F}_{c,i}^{\alpha} + {\boldsymbol F}_{a,i}^{\alpha} + \sqrt{2 \gamma / \beta} {\boldsymbol \eta}_i^{\alpha} (t)
%\label{eqn:eom}
%\end{equation}
%where $i \in \left\{1,...,N_{\rm dimers}\right\}$ indexes the dimers, $\alpha \in \left\{1,2\right\}$ indexes the monomers on dimer $i$, ${\boldsymbol p}_i^{\alpha} = m \dot{\boldsymbol x}_i^{\alpha}$ is the momentum of monomer $i$-$\alpha$, ${\boldsymbol x}_i^{\alpha}$ its position, $m$ the monomer mass, $\gamma$ a linear bath friction coefficient, ${\boldsymbol F}_{c,i}^{\alpha}$ the sum of conservative forces acting on the monomer, ${\boldsymbol F}_{a,i}^{\alpha}$ the active force acting on the monomer, $\beta \equiv 1/k_B T$ the inverse of the Boltzmann constant times the temperature, and ${\boldsymbol \eta}_i^{\alpha}(t)$ a delta-correlated Gaussian white noise.
%
%The monomers are subjected to two conservative forces: the harmonic bond linking two monomers in the same dimer, and a Weeks-Chandler-Andersen (WCA) repulsion between non-bonded monomers (WCA CITATION). The harmonic bond is characterized by a stiffness constant $k = 100 \, \varepsilon \, \sigma^{-2}$ and equilibrium length $d_0 = \sigma$, set equal to the monomer diameter $\sigma$ in all of our simulations. The WCA potential is characterized by an energy parameter $\beta \varepsilon = 1$ and the monomer diameter $\sigma$.
%
%The active forces acting on each monomer of a given dimer are of equal magnitude $F_a$ and opposite direction, ${\boldsymbol F}_{a,i}^1 = - {\boldsymbol F}_{a,i}^2 = - F_a \varepsilon \cdot \hat{\boldsymbol d}_i$, such that they exert no net force on the dimer but result in an active torque of average magnitude $F_a d_0$. The active forces are implemented in LAMMPS using code available at (CORY).
%
%The boundary current and forced run profiles are both obtained by generating an initial grid of dimers, initializing this configuration at the proper active force and imposed force for $10^7$ timesteps, and then using this initialized configuration to generate $\mathcal{O}(10^2)$ independent trajectories, each initialized using a different random seed for the Langevin thermostat. Error bars are obtained by dividing the trajectories into ten groups and block averaging.



\end{document}
%
% ****** End of file apssamp.tex ******
